% Preamble dùng chung cho main.tex (tiếng Anh) và main_vi.tex (tiếng Việt). Nạp sau fontspec và babel.
% ---- Trình bày ----
\usepackage[margin=2.5cm]{geometry}
\usepackage{microtype}
\setlength{\emergencystretch}{2em}
\usepackage{setspace}
\setstretch{1.1}
\usepackage{fancyhdr}
\pagestyle{fancy}
\fancyhf{}
\fancyfoot[C]{\thepage}
\setlength{\headheight}{15pt}

% ---- Toán, hình, bảng ----
\usepackage{amsmath,amssymb}
\usepackage{graphicx}
\usepackage{xcolor}
\usepackage{booktabs}
\usepackage{tabularx}
\usepackage{caption}
\usepackage{subcaption}
\usepackage{float}

% ---- Tài liệu tham khảo (author-year) ----
\usepackage[round]{natbib}
\bibliographystyle{plainnat}

\usepackage[colorlinks=true,linkcolor=blue!60!black,citecolor=blue!60!black,urlcolor=blue!60!black]{hyperref}

% ---- Macro cho vùng cần điền ----
% Tìm mọi chỗ còn trống: grep -rn "fillin\|fillpara\|Missing figure" reports/latex
% \fillin{gợi ý}: chỗ trống ngắn trong câu hoặc trong ô bảng.
\newcommand{\fillin}[1]{\textcolor{red}{[#1]}}
% \fillpara{gợi ý}: khung trống cho cả đoạn văn.
\newcommand{\fillpara}[1]{\par\noindent\fcolorbox{red}{red!4}{\parbox{\dimexpr\linewidth-2\fboxsep-2\fboxrule\relax}{\small\color{red}\textbf{TODO.} #1}}\par\medskip}
% Tên file, tên mô hình có dấu gạch dưới.
\DeclareRobustCommand{\code}[1]{\texttt{\detokenize{#1}}}
\newcommand{\src}[1]{\nolinkurl{#1}}

% Hình lấy từ thư mục figures/ của repo. Trên Overleaf: tải hình vào figures/ rồi sửa thành {figures}.
\newcommand{\figdir}{../../figures}
% \reportfig[tuỳ chọn includegraphics]{tên file}: chèn hình, nếu chưa có file thì hiện khung đỏ.
% Mặc định rộng bằng dòng và không cao quá 0,8 trang.
\newcommand{\reportfig}[2][width=\linewidth,height=0.8\textheight,keepaspectratio]{%
  \IfFileExists{\figdir/#2}%
    {\includegraphics[#1]{\figdir/#2}}%
    {\fcolorbox{red}{white}{\parbox[c][4cm][c]{0.9\linewidth}{\centering\color{red}Missing figure: \texttt{\detokenize{#2}}}}}}
